本组把基础技术方向和扩展模块放在一起安排。四名成员分别负责前端、后端、数据库和测试，并各自跟进看板、红包积分、AI 与高德地图等扩展内容。表\ref{tab:team}列出小组确认的主责方向和协作接口。
\begin{longtable}{L{.15\linewidth}L{.16\linewidth}L{.16\linewidth}L{.42\linewidth}}
\caption{小组实际分工}\label{tab:team}\\\toprule 成员&基础方向&扩展模块&协作接口与联调关注点\\\midrule\endfirsthead\toprule 成员&基础方向&扩展模块&协作接口与联调关注点\\\midrule\endhead
范文丽&前端&看板&对齐后端响应与页面状态，负责看板数据展示。\\
李承文&后端&红包积分&衔接订单、优惠券与积分规则，处理计价、支付和取消回退。\\
赵紫怡&数据库&AI&衔接业务数据与查询，维护真实菜单、订单授权和 AI 调用。\\
仲禹潼&测试&高德地图&组织功能与回归验证，处理定位、导航和地图失败场景。\\\bottomrule
\end{longtable}
一个需求通常需要多人一起完成。以 AI 点餐为例，候选商品来自数据库和 AI 模块，展示与确认由前端处理，加入购物车仍要经过后端规则，最后由测试走完整流程。高德地图也需要页面入口、后端数据和异常处理一起配合。

测试成员负责组织验证，每位成员复查自己提交的代码；前端主责成员也参与接口、数据和业务链路联调。小组保留技术主责，同时共同理解完整业务链路，为每项需求安排联调和验收责任。
